\section{Results and Analysis}
\label{sec:results}
\subsection{Directional task success}
Table~\ref{tab:effectiveness} contains 70,290 terminal outcomes: nine methods, two models, and 3,905 inputs per configuration. Sampling with verification achieves the highest Macro10 rate on both models (81.13\% and 86.18\%), followed by GoT (70.62\% and 83.26\%) and Reflexion (45.75\% and 53.12\%). These three procedures use property tools, but differ in generation, search, selection, and feedback. Their ordering does not isolate the contribution of any one component. LM-voted ToT, which uses the generator checkpoint to evaluate states without external property feedback, achieves 30.28\% and 21.26\%.

Rankings also depend on the backbone and task. Llama exceeds Mistral by 12.64 points for GoT and 17.78 points for four-shot ICL, but trails it by 7.88 points for s1-style budget forcing. Sampling leads all ten Mistral tasks. On Llama, GoT leads on BDPQ and BHMQ and ties Sampling on HMPQ. The BHMQ difference is one input out of 118, illustrating why per-task counts matter. Self-Refine with one revision scores 34.49\% and 31.99\%; this setting does not establish what further revisions would achieve.

\subsection{Candidate coverage and final selection}
On the unchanged width-20 Beam pools, any-success coverage is 83.38\% for Mistral and 80.76\% for Llama; returning model rank~0 gives only 35.74\% and 33.03\%. For the same recorded Sampling calls, returning slot~0 gives 32.36\% and 28.04\%. The verifier retrieves successful later candidates for 1,872 and 2,275 inputs, raising Macro10 by 48.76 and 58.14 percentage points. This comparison changes output selection while holding the recorded generation fixed. The first-candidate readout is a diagnostic on forty-candidate calls, not a separately executed one-sample experiment.

For the declared Sampling output $y^*$,
\begin{equation}
 S(x,y^*,T)=\max_{0\leq i<20}S(x,y_i,T).
\end{equation}
If any prefix candidate succeeds, the first passing candidate is returned; otherwise the fallback also fails. The equality therefore follows from the selector, including the failure case. It explains why high single-output success can coexist with much lower first-candidate success. It does not demonstrate a learned selector, better candidate generation, or independent verification of the proxy objectives. Appendix~\ref{app:readouts} reports the readouts and distinguishes them from published \gellmo{} coverage results.

\subsection{Generation work differs from output count}
Table~\ref{tab:workload} shows that a common final $K=1$ does not collapse internal computation to one generation. Exact Voting uses 400 molecular proposals; Sampling uses 40, GoT 28, and ToT 20 plus ten judge responses. Self-Refine generates two molecular proposals but four model responses. Reflexion averages 2.37 and 2.27 actor trials, reflecting its success-stopping rule rather than a fixed three-trial execution for every input.

On Mistral, Exact Voting records about 26,423 returned output tokens per input, compared with 2,596 for Sampling, yet its terminal SR is lower. This observation shows that greater returned-generation volume does not guarantee greater success across these configured procedures. It is not an ablation of voting budget: their prompts, decoding, and selection differ. Likewise, token totals across the two backbones use different vocabularies. Neither the method ranking nor these counts support an end-to-end latency or compute-efficiency claim.

\subsection{Improvement magnitude and molecular quality}
Table~\ref{tab:quality} adds information that directional SR alone does not capture. Across all ten tasks, GoT achieves magnitude-sensitive success rates of 41.27\% and 62.13\%, compared with 19.83\% and 25.60\% for Sampling. Thus Sampling more often returns a molecule with some improvement in every target, while GoT more often meets all minimum-improvement thresholds. The verifier stops at the first directional pass; it does not optimize these stronger thresholds. GoT instead uses a task-specific property-based score during search and final ranking. These differing objectives help interpret the results, but the benchmark does not isolate the score's causal effect from other pipeline differences.

GoT's final validity is 100\% on both models, while LM-voted ToT reaches 75.47\% and 50.97\%. Invalid or empty final outputs therefore form a substantial observed failure component for this ToT configuration. A valid molecule can still fail one or more property objectives, so validity alone cannot explain its entire success gap.

On their respective successful subsets, GoT yields RI$_{0.1}$ of 1.691 and 2.223, with Sim of 0.393 and 0.257. Sampling yields RI$_{0.1}$ of 0.947 and 1.153, with Sim of 0.561 and 0.453. Greater average relative improvement therefore coincides with lower source similarity in this comparison. GoT also has lower conditional SAS, but neither SAS nor the proxy-property gains establish practical synthesizability or experimental utility. Because these quality means condition on different successful inputs, they should be read alongside the full-denominator rates rather than as paired quality gains.

The three tables consequently address complementary questions: whether the returned molecule succeeds, how much generation the procedure records, and what quality its outputs exhibit. The separate reference-relative difficulty diagnostic appears in Appendix~\ref{app:reference_difficulty}; its Mistral-derived groups are not an additional model-independent ranking.
